\chapter{Balanitis}

\section*{Definition}
Balanitis is inflammation of the glans penis, often involving the foreskin (balanoposthitis). It is characterized by erythema, swelling, pain, and discharge on the glans, with recurrent episodes common in uncircumcised men.

\section*{When to See a Doctor}

\begin{itemize}
 \item Severe pain, marked swelling, discharge, fever, or difficulty urinating
 \item Recurrent episodes despite improved hygiene
 \item Associated with fever, systemic symptoms, or suspected sexually transmitted infection
 \item The foreskin is pulled back and cannot be returned over the glans (paraphimosis), or urination is obstructed; this is an emergency.
 \item A persistent ulcer, lump, bleeding area, or change in colour does not improve with treatment.
\end{itemize}

\section*{How It Develops}
Inflammation of the glans can result from infection, irritation, allergy, or an inflammatory skin disorder. Moisture and friction beneath the foreskin, diabetes, antibiotics, and immune suppression can increase the risk of Candida or bacterial infection, but infection is not present in every case.

\section*{Causes}
\begin{itemize}
 \item \textbf{Infectious:} Candida, bacterial infection, herpes simplex, and selected sexually transmitted infections
 \item \textbf{Irritant:} Soaps, detergents, lubricants, spermicides, excessive washing, chlorinated water
 \item \textbf{Dermatologic:} Contact dermatitis, lichen planus, lichen sclerosus (balanitis xerotica obliterans), psoriasis, seborrheic dermatitis, fixed drug eruption
 \item \textbf{Systemic:} Diabetes, reactive arthritis (circinate balanitis), or immune suppression
 \item \textbf{Anatomical:} Phimosis (tight foreskin trapping debris), poor hygiene
\end{itemize}

\section*{Related Symptoms}
\begin{itemize}
 \item Phimosis, paraphimosis, penile discharge, painful urination, penile rash, genital itching, meatal stenosis
\end{itemize}
\textbf{Important:} Recurrent or difficult-to-treat balanitis can justify testing for diabetes and review for an underlying skin disorder or infection.

\section*{Medical Evaluation}
\begin{itemize}
 \item Visual inspection: note erythema, swelling, discharge, fissures, phimosis, or lesions
 \item Swabs or other tests are considered when discharge, ulceration, recurrent disease, or an atypical appearance suggests infection; testing is guided by the examination.
 \item Screen for sexually transmitted infections when indicated by sexual history, lesions, or exposure.
 \item Glucose or HbA1c is considered for recurrent, severe, or otherwise unexplained balanitis.
 \item Persistent, atypical, treatment-resistant, or suspicious lesions may need dermatology/urology assessment and biopsy to exclude inflammatory disease or malignancy.
\end{itemize}

\section*{Treatment Options}
\begin{itemize}
 \item Antifungal: topical clotrimazole or miconazole for candidal balanitis
 \item Antibiotics are used only when bacterial infection is diagnosed or strongly suspected; the drug and route depend on the cause.
 \item A clinician may prescribe a low-potency topical corticosteroid for an inflammatory cause; lichen sclerosus often requires specialist-directed treatment.
 \item Recurrent disease or problematic phimosis may require urology referral; circumcision is one possible option, not an automatic treatment.
 \item Optimize glycemic control if diabetes is present
\end{itemize}

\section*{Self-Care and Home Management}
\begin{itemize}
 \item Clean the glans gently with warm water and avoid harsh soaps
 \item If the foreskin retracts comfortably, gently retract it for cleaning and drying, then return it to its normal position. Never force it.
 \item Do not start an antifungal, antibiotic, or steroid cream unless a clinician or pharmacist advises that it is appropriate; the wrong treatment can irritate the skin or delay diagnosis.
 \item Avoid sexual activity until inflammation resolves
\end{itemize}

\section*{References}
\begin{thebibliography}{99}
\bibitem{balanitis:ref1} Edwards S, Bunker CB, et al. ``Balanitis.'' \textit{BMJ}. 2019;364:l399.
\bibitem{balanitis:ref2} Wray AA, Khetarpal S, et al. ``Balanitis: diagnosis and treatment.'' \textit{Am Fam Physician}. 2021;103(10):606-612.
\bibitem{balanitis:ref3} Bunker CB. ``Male genital dermatoses: a review.'' \textit{Clin Dermatol}. 2011;29(5):507-516.
\bibitem{balanitis:ref4} Wein AJ, Kavoussi LR, et al. \textit{Campbell-Walsh Urology}, 12th ed. Elsevier, 2020.
\bibitem{balanitis:ref5} DermNet New Zealand. ``Balanitis.'' Accessed September 2026. https://dermnetnz.org/topics/balanitis.
\bibitem{balanitis:ref6} DermNet New Zealand. ``Lichen Sclerosus in Men.'' Accessed September 2026. https://dermnetnz.org/topics/lichen-sclerosus-in-men.
\end{thebibliography}
